% 来源：OEBPS/Text/01_50.xhtml（OPF spine 顺序；仅第一册）
\cleardoublepage
\BookEntry[罗杰·费希尔等《谈判力》]{罗杰·费希尔等《谈判力》}

\BookCover{images/书单33谈判力.jpg}{[美] 罗杰·费希尔、威廉·尤里、布鲁斯·巴顿著| 王燕、罗昕译| 中信出版社2012年版}

\begin{BookSummaryBox}
{\centering \textbf{哈佛大学关于谈判的最权威图书\\  20年来最为经典的谈判图书}\par}
\end{BookSummaryBox}

\subsubsection*{这本书谈什么？}

在日常生活中，每个人都不可避免会成为一个谈判者。虽然每天要与谈判相伴，但并不是每个人都能谈出好的结果。在常见的谈判策略中，人们使用的大多属于温和或强硬的方式。温和的方式总是尽量避免摩擦和冲突，为达成共识，往往很快让步；而强硬的方式则以战胜对方为目标，坚守立场，但常会两败俱伤。
\par

在这本书里，作者基于“哈佛大学谈判项目”提出了第三种谈判方式，即“原则谈判方式”。这种方式根据事情本身的是非曲直寻求解决方案，强调把人和事分开，着眼于利益而非立场，当双方利益发生冲突时，让谈判结果基于某些公平的标准，而不是以某一方意志为转移。
\par

尽管谈判的形式与内容千差万别，但原则谈判方式在各种谈判中都完全适用。
\par

\subsubsection*{芒格为什么谈起这本书？}

芒格在“专业人士需要更多的跨学科技能”的演讲中，鼓励大家多多使用跨学科方法，多多实行“拿来主义”。因为在软科学界研究发现，有一种研究方法非常好，就是“拿来主义”法，不论看中其他学科什么知识，只管拿过来用，以免在单一学科中故步自封。
\par

谈到具体的例子，他说道：
\par

\begin{quotation}
只要多多实行“拿来主义”，许多软科学就能减少“铁锤人倾向”引起的错误。例如，在我们的同学罗杰·费希尔的领导下，许多法学院将其他学科的成果应用到谈判研究中去。罗杰那本充满智慧和道德感的谈判著作到目前已经卖掉300多万册，他可能是我们班最有成就的同学。这些法学院还吸取了大量有用的经济学知识，甚至用博弈论来更好地理解经济竞争的原理，从而制定出更有针对性的反垄断法。
\par

罗杰·费希尔的《谈判力》，就是“拿来主义”的典范性佳作。
\par
\end{quotation}

\subsubsection*{《谈判力》作者罗杰·费希尔等人如是说}

“如何避免陷入攻击和辩解的恶性循环中？答案是不要回击。假如对方直截了当地表明了立场，不要拒绝，而是透过立场看利益，把对方的立场当作一种可能的选择，找到其遵循的原则并考虑改进的方法；当对方反驳你的观点时，不要辩解，欢迎批评和建议，把批评引向建设性轨道的另一种方法是，转换双方的位置，征求对方的意见，问问他们如果站在你的角度他们会怎么做；当他们人身攻击，你应不动声色，让对方把怒气发出来，倾听对方的诉说，表现出你愿意尊重他们的意见。”
\par

“要想在纷繁复杂的人际问题中找到出路，最好从以下三个基本方面着手：认知、情绪和交流。各种人际关系问题都能从这三个方面找到症结所在。”
\par

“谈判的任何办法都可以通过三方面的标准来衡量：如果有达成共识的可能，就应该达成明智的协议；谈判必须有效率；增进或至少不损害双方的关系（明智的协议是指协议尽可能保障双方的合法利益，公平解决双方的利益冲突，协议持久性强，并考虑了社会效益）。”
\par
